\subsection{The two-level metric bridge and the remaining weighted height bound}
\label{sec:geometry-two-level-weighted-bridge}

We record the precise arithmetic intersection consequence of the local
calculations. This separates the geometric linkage of the two original
models from the subsequent estimates of their weight polynomials.

\begin{lemma}[Vertical comparison on the actual common model]
\label{lem:geometry-two-level-vertical-bridge}
Suppose the two cutoffs \(B_0\le B_1\) have a common normal
projective comparison model with their specified root line bundles
\(X_i,Y_i\), pulled back with the original metrics. Assume:
\begin{enumerate}
\item their generic identifications are isometries, and the
  differences of the finite lattices are supported on the new
  correction primes \(S(B_1)\setminus S(B_0)\);
\item on each such fibre the new roots in each prescribed label
  are the graded lines of the trivial flagged bundle, as in (HF1);
\item the old and new root bundles at that prime obey the actual
  two relative Chern identities (RC4)--(RC5), with their
  original labelling and dual conventions.
\end{enumerate}
Then, for every \(k\ge1\) and every complementary root monomial
\(Q\) of degree \(d+1-k\), the difference of the original arithmetic
intersections of \(p_{k,\tau}Q\) is exactly the sum of their local
relative Chern comparisons with \(Q\) evaluated on the old roots,
each multiplied by the prime's logarithm and the original
arithmetic normalization factor. There is no additional
mixed-root remainder.
\end{lemma}
\begin{proof}
Use arithmetic first Chern classes \(x_i,y_i\) of the two
metrized line bundles, and the equality
\[
 p_{k,\tau}(x)Q(x)-p_{k,\tau}(y)Q(y)
 = [p_{k,\tau}(x)-p_{k,\tau}(y)]Q(x)
   +p_{k,\tau}(y)[Q(x)-Q(y)] .
\]
The generic isometries identify \(x_i-y_i\) with arithmetic
classes of vertical Cartier divisors, with zero archimedean
part. Telescope each monomial in \(Q(x)-Q(y)\) using these
differences. The degree of each resulting term in the last
summand is zero: intersect on its supporting special fibre,
where \(p_{k,\tau}(y)=0\) by the exact trivial flag and the
Whitney formula. This remains true on a nonreduced special
fibre, by applying the intersection to its component cycles
with their multiplicities. The first summand is the arithmetic
realization of the relative Chern comparison of the two
root bundles; this is precisely the power sum of the vertical
line differences multiplied by their mixed-degree polynomial
and by \(Q(x)\). Its evaluation is the local number in (RC4),
with \(\log p\) and the original residue degree multiplicities.
The projection formula shows invariance under a common
birational refinement and, for normalized arithmetic degrees,
under a finite common coefficient extension. For the latter,
both the total arithmetic degrees and the field-degree
denominator are multiplied by the extension degree.
Linearity gives the same identity for complementary
polynomials. This proves the stated comparison, without
discarding a principal vertical divisor before its fibre
intersection is taken.
\end{proof}

This lemma gives the usual two-level power-sum comparison
once the same-model and isometry hypotheses have been supplied.
It does not obtain those hypotheses by changing the metrics
or by choosing unrelated root bundles at the two cutoffs.
In particular the first-order local argument in
Section~\ref{sec:geometry-witt-additive-lattice} must still
be matched to the global charts as stated there.

Here is the exact numerical consequence of the already proved
base test for the original weighted-height reduction.
Use the notation of \texttt{tower-intersection-reduction.md}:
\(\Lambda_B=\log(B\operatorname{rd}(E))+|\lambda_E|\),
\(f=|\Psi|\), \(N=p+q\), \(3q+1\le N\), \(f\ge2\),
and \(L(u)=\sum_{\tau,i}|u_{\tau i}|\),
\(W(u)=\sum_\tau(u_{\tau0}-u_{\tau,N-1})\).
These \(p,q\) denote the signature, not a correction prime.
Suppose the exact two-polynomial comparison (2.1) of that
note and its bounds (2.4) have been established on the
same chosen tower, with the original common signed
determinant coefficient \(|c_B|\le4\Lambda_B\).
The signed base-test theorem now supplies
\(\beta_B\ge-107\Lambda_B/2\).
Consequently the block-sensitive elimination gives
\[
 \sigma_B(u)\ge
 -K_BL(u)-2qW(u),
 \quad
 Z=\tfrac{107}{2}\Lambda_B+
       \tfrac12\log(N+1)+\frac8{5fq},
 \quad
 K_B=\max\{Z,(q/p)Z+(N/p)4\Lambda_B\}.
\]
For the original \(B\ge\max(128,4fN^2)\) one has
\(Z\le55\Lambda_B\): indeed
\(\log(N+1)\le\log B\), \(8/(5fq)\le4/5\le\log B\).
As \(q/p\le1/2\), \(N/p\le3/2\), the other entry in the
maximum is at most \((55/2+6)\Lambda_B\).
Thus
\begin{equation}\label{eq:geometry-conditional-mean55}
 \sigma_B(u)\ge-55\Lambda_BL(u)-2qW(u).
\end{equation}
With equal component volumes and the actual prescribed
norm/count-volume comparison \(s_Z=\sigma_Z-2c_2(E,B)L(u)\)
of \texttt{mean-volume-supplier.md}, the finite-place
bound \(c_2\le40(\log B+\log\operatorname{rd}(E))\) yields
\[
 \frac1{|\mathcal Z|}\sum_Zs_Z\ge
       -135\Lambda_BL(u)-2qW(u).
\]
All constants are independent of the growing datum.
The displayed weighted bound is conditional only on the
specified remaining same-tower comparison and its flag
height/remainder bounds and norm dictionary. Neither
the source's name of a comparison theorem nor the
local test identities alone establishes these inputs.
The finite section-rank comparison and the final
same-window arithmetic budget are separate.
